\documentclass{article}

\usepackage{microtype}
\usepackage{graphicx}
\usepackage{booktabs}
\usepackage{amsmath}
\usepackage{amssymb}
\usepackage{hyperref}
\newcommand{\theHalgorithm}{\arabic{algorithm}}
\usepackage{icml2026}

\icmltitlerunning{Auditable Branching World Models for Forecasting}

\begin{document}

\twocolumn[
\icmltitle{WorldFork: Auditable Branching Socio-Institutional\\World Models for LLM Forecasting Agents}

\begin{icmlauthorlist}
\icmlauthor{Anonymous Authors}{anon}
\end{icmlauthorlist}
\icmlaffiliation{anon}{Anonymous Institution}
\icmlcorrespondingauthor{Anonymous Authors}{anonymous@example.com}
\icmlkeywords{forecasting, agents, world models, evaluation, uncertainty}

\vskip 0.3in
]

\printAffiliationsAndNotice{}

\begin{abstract}
LLM forecasting systems usually return a probability and rationale, but they rarely expose the counterfactual rollouts, endpoint evidence, social assumptions, and unresolved uncertainty behind that forecast. We present WorldFork, a forecasting protocol that turns a scenario into checkpointed timelines with branch lineage, endpoint ledgers, path mass, report provenance, and explicit socio-institutional state. We build an evaluation package with 108 public initialization and audit cards, a 24-card retrospective resolved forecasting pilot, and long-horizon branch-stress tasks for lineage and report grounding. In the current run, static leakage checks pass, source fetch verification succeeds for 39 of 40 resolution URLs, GPT-5.4 and DeepSeek v4 Flash direct LLM baselines complete on all 24 resolved cards, and a deadline-aware E3 branching core-12 run completes. The canonical E3 WorldFork score aggregates explicit yes/no endpoint path mass across branches. A direct-prior sensitivity analysis reuses existing E2 calls and shows that blending the DeepSeek structured direct prior with E3 branch mass improves the core-12 in-sample score, though leave-one-out tuning is not yet robust. Tick counts are caps, not targets: endpoint-ledger and path-mass resolution is the natural stopping criterion, while unresolved or insufficient mass resumes the same Big Bang to preserve lineage and cost accounting.
\end{abstract}

\section{Introduction}

Forecasting is increasingly used to test whether general-purpose AI systems can reason under uncertainty about the real world. A useful forecast, however, is more than a scalar probability: proper scoring rules measure accuracy \citep{brier1950verification,gneiting2007strictly}, while forecasting practice also depends on exposing assumptions and uncertainty \citep{tetlock2015superforecasting}. Institutional and social events depend on actors, authorities, incentives, trust and dependency relationships, procedural deadlines, information flows, and unresolved evidence. Single-shot LLM forecasts can mention these factors in prose, but their internal assumptions are difficult to audit or compare after the fact.

WorldFork addresses this gap by representing a forecast as an auditable branching world model. A scenario is initialized into actors, cohorts, candidate endpoints, graph layers, and initial socio-institutional state. The runtime advances timelines in ticks, forks alternate continuations when mechanisms diverge, tracks endpoint evidence in ledgers, records path mass, and generates reports with provenance. The purpose is not to claim that branching simulation automatically improves probabilistic accuracy. The purpose is to expose the forecast object: which assumptions led to which endpoint mass, where uncertainty remained, and which report claims are grounded in logged artifacts.

This paper makes three contributions. First, it specifies a branching forecast protocol with checkpointed Big Bangs, ticks, branch lineage, endpoint ledgers, path mass, and report provenance. Second, it treats socio-institutional state as audit data: actors, cohorts, trust/dependency/conflict/influence/coalition/exposure graphs, sociology signals, source packets, and bounded emotion observations. Third, it provides a pilot benchmark and evaluation workflow covering initialization quality, resolved forecast scoring, long-horizon audit metrics, and social-state consistency.

\section{Method}

\textbf{Initialization.} A public scenario card is converted into a Big Bang containing typed actors and cohorts, scenario constraints, candidate endpoints, graph layers, and initial social state. For resolved forecast cards, endpoints are binary yes/no outcomes. For dossier and calibration cards, endpoints include institutional resolutions plus unresolved or abstention states. Initialization is scored for schema completeness, actor recall, authority fidelity, constraint preservation, endpoint extraction, graph quality, prompt-injection resistance, and uncertainty honesty.

\textbf{Ticks and branches.} A multiverse advances in ticks. Each tick updates actor positions, social signals, events, endpoint evidence, and logs. Tick counts are maximum budgets rather than required stopping points: a run may stop when endpoint-ledger state and path-mass artifacts show resolution, while unresolved or insufficient mass is handled by continuing the existing Big Bang rather than reinitializing it. A branch policy controls maximum depth, active multiverses, branches per tick, and branch-score threshold. Short resolved runs compare no-branch and branching policies. Long-horizon runs stress whether lineage, branch locality, path mass, and endpoint ledgers remain inspectable over many ticks.

\textbf{Endpoint ledgers and reports.} Endpoint ledgers track whether candidate endpoints are supported, contradicted, unresolved, or not yet checkable. Reports are scored by sampling claims and verifying support from ticks, ledgers, logs, or source packets. For resolved scoring, yes/no probability is normalized over yes and no when unresolved mass is present; unresolved mass is still reported separately. A forecast row is frozen only when terminal job artifacts and path-mass evidence support the ledger state; otherwise continuation preserves branch lineage, report provenance, and cost history.

\textbf{Socio-institutional state.} WorldFork models social assumptions as audit state, not validated psychology. Graph layers and sociology signals are useful only when they make branch divergence or report claims traceable. Emotion observations are bounded logged variables and should not be interpreted as hidden human affect.

\begin{figure}[t]
\centering
\fbox{\begin{minipage}{0.94\columnwidth}
\centering
\scriptsize
Scenario card $\rightarrow$ Big Bang initialization\\
$\downarrow$\\
Actors, endpoints, constraints, graph/social state\\
$\downarrow$\\
Ticked multiverses $\rightarrow$ branch policy $\rightarrow$ alternate timelines\\
$\downarrow$\\
Endpoint ledgers + path mass + logs\\
$\downarrow$\\
Final forecast/report with provenance
\end{minipage}}
\caption{WorldFork represents a forecast as a set of inspectable runtime artifacts rather than only a probability and rationale.}
\label{fig:protocol}
\end{figure}

\section{Benchmark and Metrics}

The benchmark contains 108 public cards: 72 existing synthetic WorldFork stress prompts and 36 added forecasting cards. The add-on contains 24 retrospective resolved forecast cards, 8 longform source-packet dossiers, and 4 adversarial/calibration cards. Private resolutions and gold checklists are withheld from forecast-producing systems until forecasts are frozen.

\begin{table}[t]
\centering
\footnotesize
\begin{tabular}{p{0.48\columnwidth}rp{0.28\columnwidth}}
\toprule
Group & Count & Primary use \\
\midrule
Synthetic stress cards & 72 & initialization and audit \\
Resolved forecast cards & 24 & Brier/log scoring \\
Longform dossiers & 8 & report and ledger grounding \\
Adversarial calibration & 4 & uncertainty and injection tests \\
\bottomrule
\end{tabular}
\caption{Benchmark composition for the ICML forecasting package.}
\label{tab:composition}
\end{table}

Forecast accuracy on resolved cards uses binary Brier score $(p_{\mathrm{yes}}-y)^2$ and clamped negative log score. WorldFork outputs may assign mass to unresolved states; primary yes/no scoring normalizes over yes/no when possible and reports unresolved mass separately. Audit runs use 0--4 rubric scores for lineage integrity, branch locality, endpoint coverage, endpoint honesty, path-mass consistency, report grounding, failure observability, cost transparency, social-state consistency, and emotion observability.

\section{Results}

\textbf{Card QA.} Static package QA passes. The run contains 108 public cards, 36 private evaluation rows, matching public/private IDs, no private resolution fields in public cards, 24/24 resolved cards with at least one resolution source, and balanced private labels (12 yes, 12 no). Automated source fetching plus browser follow-up checked 40 resolution URLs: 39 verified successfully, and the remaining gated media URL belongs to a case that already has a separate primary court-source row marked ok.

\textbf{Direct baselines.} We ran two retrieval-free OpenAI Codex \texttt{gpt-5.4} direct baselines on all 24 resolved cards. The unstructured direct prompt achieved mean Brier 0.242492 and mean log score 0.701698. The structured direct prompt achieved mean Brier 0.238363 and mean log score 0.677498. Both baselines have zero unresolved mass because the output contract requires a normalized yes/no distribution.

\begin{table}[t]
\centering
\footnotesize
\begin{tabular}{lrrrr}
\toprule
Condition & $n$ & Brier & Log & Unres. \\
\midrule
Direct core12 & 12 & 0.2199 & 0.6733 & 0.0000 \\
Structured core12 & 12 & 0.2155 & 0.6371 & 0.0000 \\
DeepSeek structured core12 & 12 & 0.1950 & 0.5918 & 0.0000 \\
WF branching core12 & 12 & 0.2246 & 0.6133 & 0.0000 \\
WF + DeepSeek structured 50/50 & 12 & 0.1906 & 0.5534 & 0.0000 \\
\bottomrule
\end{tabular}
\caption{Resolved-card forecast scores on the 12 IDs covered by the E3 branching run. The WorldFork row aggregates explicit yes/no endpoint path mass across multiverses; the blend row reuses existing DeepSeek structured direct calls as a fixed prior with no new simulation or LLM calls.}
\label{tab:forecast}
\end{table}

\textbf{Live WorldFork evidence.} The full 108-card suite now has live initialization evidence. Initialization on one resolved card completed in 146.20 seconds with 9 actors, 9 traits, graph state, sociology baseline, and emotion baseline present. Queued initializer batches then completed the remaining public cards: the largest add-on batch reached p1's configured concurrency of 8 and drained 22 jobs in 464.64 seconds wall time, while the existing 72-card batch completed in 1796.65 seconds with p1 saturated at 8 for most of the run. The enriched initialization table reports 108/108 succeeded rows, mean actor count 8.82, mean graph edge count 34.21, and an automated evidence-grounded 0--4 rubric proxy with mean score 3.884, 108/108 pass rows, and no critical failures. This rubric proxy is a structural audit score, not a human semantic adjudication.

The main WorldFork E3 comparison uses the same deadline-aware branching core-12 runs and scores the candidate endpoint path-mass distribution across all multiverses. A single-path proxy that selects the multiverse with the largest stored path probability is retained only as a diagnostic ablation; it has mean Brier 0.4167, mean log score 1.9247, and mean unresolved mass 0.0000. The branching aggregate has mean Brier 0.2246, mean log score 0.6133, and mean unresolved mass 0.0000.

On the same 12 card IDs, GPT-5.4 direct prompting has mean Brier 0.2199 and log score 0.6733, while GPT-5.4 structured direct prompting has mean Brier 0.2155 and log score 0.6371. DeepSeek v4 Flash structured direct is stronger at Brier 0.1950 and log score 0.5918. A 50/50 blend of DeepSeek structured direct with the E3 branching aggregate scores Brier 0.1906 and log score 0.5534. The best in-sample Brier blend uses alpha=0.70 and reaches Brier 0.1877, but leave-one-out alpha tuning degrades to Brier 0.2328, so tuned alpha is diagnostic rather than a robust paper claim.

\textbf{Queue tuning.} The synchronous smoke left Celery queues idle. The queued smoke and scored validations showed active p1 \texttt{worldfork.execute\_job} tasks while p0/p2/p3 were idle. The queued initializer batches showed that p1 can run independent cases concurrently up to its configured concurrency of 8, and the three-case E3 no-branch batch confirmed that independent run jobs overlap in one shared wall-clock window. Thus Celery is functioning; queueing improves throughput across cases, while one case remains a long p1 job. The first branching validation shows that branch policy, not queue capacity, is now the main runtime multiplier.

\section{Limitations}

The resolved cards are retrospective and use partial masking, so leakage cannot be ruled out. The synthetic and dossier cards test auditability and stress behavior rather than direct real-world forecasting accuracy. Social-state metrics are rubric-based and require evidence sampling. Emotion observations are audit variables, not validated human psychology. The main WorldFork forecast comparison covers only a core-12 fallback subset, and the single-path no-branch proxy is derived from branching runs rather than a separately rerun no-branch condition. Tick caps should not be read as equal-length rollouts; they are budgets conditioned on ledger resolution and continuation state. Stronger accuracy or auditability claims require full-card branching and the pending long-horizon and social-state scoring blocks.

\section{Conclusion}

WorldFork is best viewed as an auditable forecasting substrate. Its value is the production of inspectable forecast artifacts: branch lineage, endpoint ledgers, path mass, report provenance, unresolved uncertainty, and social assumptions that can be checked after the fact. The current evidence supports the feasibility and scoring setup, while the final accuracy and auditability claims should be made only after the full WorldFork conditions are completed.

\bibliography{refs}
\bibliographystyle{icml2026}

\end{document}
